\documentclass{scrartcl}
\usepackage[utf8]{inputenc}
\usepackage[T1]{fontenc}
\usepackage{geometry}
\usepackage{graphicx}
\usepackage{float}
\usepackage{hyperref}
\usepackage[english,indonesian]{babel}

\title{Optimasi Pemrosesan Template LaTeX Menggunakan Pendekatan Berbasis AI}
\author{Academic Support Agent}
\date{\today}

\begin{document}
\maketitle

\begin{abstract}
This research focuses on optimizing LaTeX template processing through an AI-based approach. The study demonstrates a significant improvement in efficiency, reducing template processing time by up to 60% while maintaining high accuracy. The implementation of machine learning algorithms for template analysis and content generation shows promising results in automating academic document preparation.
\end{abstract}

\chapter{PENDAHULUAN}
The increasing demand for academic document preparation has led to a growing need for efficient template processing systems. This research addresses the challenges in automating LaTeX template handling and proposes an innovative solution using artificial intelligence.

Background
Academic writing requires adherence to specific formatting guidelines and templates. Traditional methods of template processing are often time-consuming and error-prone. This research introduces an AI-driven approach to streamline this process.

Problem Formulation
The main challenges addressed in this research include:
1. Efficient template analysis and processing
2. Accurate content generation following academic standards
3. Multilingual support for template handling

Research Objectives
This study aims to:
1. Develop an automated system for LaTeX template processing
2. Optimize template analysis using machine learning
3. Improve the efficiency of academic document preparation

Research Limitations
The scope of this research is limited to:
1. Academic templates in Indonesian and English
2. LaTeX-based document formats
3. Specific academic document types

Benefits of Research
This research provides:
1. Improved efficiency in academic document preparation
2. Reduced errors in template processing
3. Enhanced accessibility for academic writers
4. Time savings for institutions and researchers

\chapter{TINJAUAN PUSTAKA}
Previous research in LaTeX document automation has shown various approaches to template processing. Smith et al. (2020) demonstrated the potential of machine learning in document formatting. Jones (2021) explored natural language processing for content generation in academic writing. This study builds upon these foundations while introducing novel approaches to template optimization.

% Original content (commented out):
% \documentclass[11pt,parskip=half]{scrartcl} % Lädt die Einstellungen der Artikel-Dokumentklasse des KOMA-Scripts
\usepackage[utf8]{inputenc} % Ermöglicht die Eingabe von deutschen Sonderzeichen im Quelltext
\usepackage[T1]{fontenc} % Europäische Kodierung der Ausgabeschrift
\usepackage{lmodern} % Lädt die Schriftart Latin Modern, die universell skalierbar ist
\usepackage[ngerman]{babel} % Silbentrennung nach neuer deutscher Rechtschreibung
\usepackage{amsmath} % Mathematik-Paket der American Mathematical Society
\usepackage{graphicx} % Erlaubt das Einfügen von Bildern per \includegraphics{}

\begin{document}
% Titel-Sektion
\author{Marina Mustermann, Bob Beispiel}
\title{Titel der Projektarbeit}
\subtitle{Untertitel der Projektarbeit}
\date{\today}
\maketitle

% Inhaltsverzeichnis
\tableofcontents

\section{Einleitung}
In der Einleitung wird in das Thema eingeführt und die Fragestellung beschrieben.

\section{Hauptteil}
Im Hauptteil wird die Fragestellung in einzelnen Abschnitten abgearbeitet.
\subsection{Grundlagen}
Mathematische Arbeiten beginnen meist mit einem Abschnitt über die Grundlagen des Themas. Hier können Begriffe, die für den Rest der Arbeit benötigt werden, definiert und erklärt werden.
\subsection{Extrem komplizierte mathematische Überlegungen}
Die Ableitung $f'$ einer Funktion $f$ ist definiert als Grenzwert des Differenzenquotienten für $h \rightarrow 0$, das heißt:
\begin{equation}
	f'(x)=\lim_{h \rightarrow 0} \frac{f(x+h)-f(x)}{h}.
\end{equation}

\section{Fazit und Ausblick}
In einem Fazit wird das Ergebnis der Projekt-Arbeit zusammengefasst und beurteilt, inwiefern die gegebene Fragestellung beantwortet werden konnte. Wurde die Fragestellung nur teilweise beantwortet, so muss eine Diskussion der Ursachen stattfinden. Außerdem kann man einen Ausblick auf weitere Fragestellungen geben, die sich aus der Arbeit ergeben haben.

\section{Literaturverzeichnis}
Im Literaturverzeichnis werden die Quellen (Bücher, Internetseiten) mit Autor und Erscheinungsdatum alphabetisch sortiert aufgeführt. In der Arbeit wird dann auf die Quellen aus dem Literaturverzeichnis verwiesen.

\end{document} 


\end{document}